\subsection{Excluding a point of order $55$}
\begin{lemma}\label{lem:decic55-source-injectivity}
Reduction at $2$ is injective on $J_1(55)(\mathbb Q)_{\mathrm{tors}}$.
\end{lemma}
\begin{proof}
Put $X=X_1(55)$, of genus $81$, and let $\pi:X\to Y$ be the quotient by
$\langle\langle16\rangle\rangle$, of order $5$.
This group acts freely, including on the special fibre at $2$:
there is no nontrivial elliptic automorphism of order $5$, and
primitive cusp columns give free orbits of length $5$.
Thus $Y$ has good reduction at $2$ and genus $17$.
By \cite[Theorem~4.13 and Table~2]{DEHMZ},
\[
 J_X(\mathbb Q)[2^\infty]=C_2\times C_{16}.
\]
Since $\pi_*\pi^*=[5]$, pullback injects the rational $2$-primary
torsion of $J_Y$ into this group.

We recall the degree-independent finite calculation on $Y$.
The map $Y\to X_0(55)$ is an etale cyclic cover of degree $4$.
In the order of cusp widths $55,11,5,1$, its cusp Frobenius
elements at $2$ are $0,1,2,3$ in $C_4$.
The normalized base has genus $5$ and Picard group
$C_5\times C_{20}\times\mathbb Z$.
There is a unique character to $C_4$ with these cusp values.
The corresponding explicit quartic is identified with the modular
cover by its exact Witt equations, inverse maps and Frobenius
values on generators of the Picard group modulo $4$.
Its genus is $17$, and a split rational cusp gives the exact
constant field.

The exact finite Jacobian and the subgroup generated by normalized
finite cusp classes are
\[
 J_Y(\mathbb F_2)=[5,5,10,1200],\qquad H=[5,10,1200].
\]
Every finite cusp is the reduction of a rational effective global
cuspidal divisor of the same degree: the global and finite cusp
orbits coincide, with signature $1^4,2^2,4^2$.
Manin--Drinfeld therefore supplies rational torsion lifts of the
generators of $H$.
The $2$-primary part of $H$ has order $32$, attaining the upper
bound from $J_X(\mathbb Q)[2^\infty]$.
Consequently $J_Y(\mathbb Q)[2^\infty]$ has order $32$, and
reduction is injective on it.

The pullback $\pi^*$ is now an isomorphism on rational
$2$-primary torsion.
If $\pi^*Q$ reduces to zero, pushforward gives
$5\overline Q=0$.
Since $\overline Q$ is $2$-primary, it is zero, whence $Q=0$.
Thus reduction is injective on $J_X(\mathbb Q)[2^\infty]$.
It is always injective on odd torsion at $2$, proving the lemma.
\end{proof}

\begin{proposition}\label{prop:decic-no55}
An elliptic curve over a number field of degree $10$ has no point
of order $55$.
\end{proposition}
\begin{proof}
Put $X=X_1(55)$, of genus $81$.
Its rational Jacobian has rank zero by
\cite[Theorem~3.1]{DEHMZ}.
The projective index is $1440$, and \cite[(2.1)]{NV} gives
\[
 \operatorname{gon}_{\mathbb Q}X\geq
 \frac{325\cdot1440}{32768}=\frac{14625}{1024}>10.
\]
Suppose that a torsion point gives a rational effective divisor
$D$ of degree $10$, with entirely noncuspidal support.

The global and finite cusp orbits at $2$ coincide, with signature
\[
 1^{20},4^5,10^2,20.
\]
In particular, both degree-$10$ residue cusps are globally intact.
At denominator $11$ their field is $\mathbb Q(\zeta_{11})$,
and $2$ has order $10$ modulo $11$.
At denominator $55$ the field is
$\mathbb Q(\zeta_{55})^+$, and the projective order of $2$
modulo $55$ is $20$.
The complete primitive-column calculation also checks every
rational and quartic cusp.
If $\overline D$ is entirely cuspidal, lift it to a rational
cuspidal divisor $C$ of degree $10$.
The torsion class $[D-C]$ reduces to zero.
Lemma~\ref{lem:decic55-source-injectivity} gives $D\sim C$,
and the gonality bound forces $D=C$, a contradiction.

The complete noncuspidal closed-point counts through degree $10$
are
\[
 (0,0,0,0,0,0,20,0,40,100).
\]
An independent census checks all ordinary elliptic curves and
their quadratic twists; the supersingular trace possibilities
give no additional point of order $55$.
In degrees $7,9,10$, the numbers of qualifying elliptic curves
are $7,18,50$, respectively.
In every case the rational $55$-torsion subgroup is cyclic of
order $55$. Its $20$ generator classes form one diamond orbit.
The closed points therefore form one, two and five source
diamond orbits in degrees $7,9,10$.
A reduction with noncuspidal support consists of one such place,
with a complement of respectively three, one or zero cusps.
These complements are supported on rational cusps.

The operator
\[
 A_3=T_3-1-3\langle3\rangle
\]
annihilates $J_X(\mathbb Q)$ by Eichler--Shimura and injectivity
of torsion reduction at the good odd prime $3$.
For a rational cusp $c$, the cusp Hecke formula gives
$T_3c=c+3\langle3\rangle c$ as a literal divisor.
Thus the proposed divisor makes the $A_3$-image of its
noncuspidal reduction principal.

The accepted degree-independent source calculations on the
septic orbit and both nonic orbits give nonprincipal $A_3$-images.
The added rational cusp complements contribute zero, so these
calculations apply unchanged here.
For the five new degree-$10$ representatives, fresh exact
Riemann--Roch calculations give
\[
 h^0(T_3P-P-3\langle3\rangle P)=0
\]
in every case.
Four representatives have two Frobenius orbits of cyclic
$3$-kernels, of lengths $1$ and $3$, with image degrees $10$ and
$30$ and norm weights one.
The remaining representative has four rational kernels and four
degree-$10$ images, again with norm weights one.
Independent arithmetic reconstructs all $20$ marked cyclic
$3$-isogenies and the complete $3$-division polynomials, and checks
the Tate coordinates and every norm weight.
The prime certificates use two exact ideal generators in a
single absolute model over $\mathbb F_2(x)$.
A second native calculation rebuilds all $21$ distinct prime
ideals, verifies their marked coordinates and every divisor norm,
and reproduces all five zero Riemann--Roch dimensions.
Since Hecke operators commute with diamonds, these five
nonprincipality tests exclude all $100$ degree-$10$ source
places. This completes the proof.
\end{proof}
